\chapter{Całka Lebesgue'a - własności}
\label{chap:calka-lebesguea-wlasnosci}

\section{Własności całki funkcji całkowalnej}
\label{sec:wlasnosci-calki-funkcji-calkowalnej}

\begin{twierdzenie}
\label{tw:wlasnosci-calki-funkcji-calkowalnej}
Niech $E\in\mathcal F$ i niech $f$ oraz $g$ będą całkowalne na $E$.
\begin{enumerate}
  \item Dla $c\in\R$ funkcja $cf$ jest całkowalna na $E$ oraz
        $\displaystyle\int_E cf\,d\mu=c\int_E f\,d\mu$.
  \item Funkcja $f+g$ jest całkowalna na $E$ oraz
        $\displaystyle\int_E(f+g)\,d\mu=\int_E f\,d\mu+\int_E g\,d\mu$.
  \item Jeżeli $f\leq g$ na $E$, to
        $\displaystyle\int_E f\,d\mu\leq\int_E g\,d\mu$.
  \item Jeżeli $F\in\mathcal F$, $E\cap F=\varnothing$ i $f$ jest
        całkowalna na $E\cup F$, to
        \[
          \int_{E\cup F}f\,d\mu=\int_E f\,d\mu+\int_F f\,d\mu.
        \]
\end{enumerate}
W punkcie 2) dla funkcji o wartościach rozszerzonych sumę można
określić jako zero w punktach, w których pojawia się
$(+\infty)+(-\infty)$ lub $(-\infty)+(+\infty)$.
\end{twierdzenie}

\begin{proof}[Dowód punktu 2)]
Wystarczy rozpatrzyć $E=X$, ponieważ całki po $E$ są całkami funkcji
$f\mathbf{1}_E$ i $g\mathbf{1}_E$ po $X$.

Funkcja całkowalna $h$ jest skończona prawie wszędzie. Istotnie, dla
$N=\{|h|=+\infty\}$ zachodzi $n\mathbf{1}_N\leq|h|$ i stąd
$n\mu(N)\leq\int_X|h|\,d\mu<+\infty$ dla każdego $n\geq1$.
Dla $f$ i $g$ możemy więc wykonać rachunek po wyzerowaniu obu funkcji
na sumie ich zbiorów wartości nieskończonych, bez zmiany całek części
dodatnich i ujemnych (uwaga~\ref{uw:calka-a-zbior-miary-zero}).

Z nierówności $|f+g|\leq |f|+|g|$ i własności całki funkcji
nieujemnej dostajemy
\[
  \int_X|f+g|\,d\mu
  \leq\int_X|f|\,d\mu+\int_X|g|\,d\mu<+\infty.
\]
\begin{samepage}
Zatem $f+g$ jest całkowalna. Ponieważ nierówność $(f+g)^+\leq f^++g^+$
może być ostra, korzystamy z rozkładu
$(f+g)^+-(f+g)^-=f^+-f^-+g^+-g^-$. Po przeniesieniu części ujemnych
otrzymujemy tożsamość funkcji nieujemnych
\[
  (f+g)^+ + f^- + g^- = f^+ + g^+ + (f+g)^-.
\]
Stosując addytywność całki po obu stronach, otrzymujemy
\begin{align*}
  \int_X(f+g)^+\,d\mu+\int_X f^-\,d\mu+\int_X g^-\,d\mu
  ={}&\int_X f^+\,d\mu+\int_X g^+\,d\mu\\
    &+\int_X(f+g)^-\,d\mu.
\end{align*}
Wszystkie składniki są skończone. Po przeniesieniu całek części
ujemnych dostajemy $\int_X(f+g)\,d\mu=\int_X f\,d\mu+\int_X g\,d\mu$.
\end{samepage}
\end{proof}
